\paragraph{The integral maps at one cyclotomic level.}
\label{par:row-metric}
\hypertarget{source-row-metric}{}
\begin{paramdefinition}[Local field, signed intervals and integral map]\label{hyp:first-level-row}
Write \(p=\ell\), let \(K/\mathbb Q_p\) be the local field, and let
\(K_0\) be its maximal unramified subfield. Its residue degree,
ramification index and absolute degree are
\[
 n=[K_0:\mathbb Q_p],\qquad e=[K:K_0],\qquad D=[K:\mathbb Q_p]=ne.
\] In the application \(e=2\).
The calculation also applies at \(e=1\).
A first-level row has odd residue characteristic,
\[
 K\subset K_0(\mu_p),
\]
and a coefficient
DVR \(O\) containing the images of the rings of integers
under all embeddings of \(K_0(\mu_p)\). All the following lengths are
divided by \({\rm length}_O(O/pO)\). Let \(F:T_1\twoheadrightarrow T_2\)
be the given integral crystalline surjection, with \(r={\rm rank}_O T_2\).
Write \(h,S,\eta\) for the integral height, width and central radius.
If the sign at \(\tau\) is \(d_\tau\in\{+1,-1\}\), the corresponding outer jumps \(a\) satisfy
\[
 b_\tau\le d_\tau a\le b_\tau+S.
\]

The signed jump data include a central interval \([-\eta,\eta]\) and a bound \(m\) on the number of central target jumps at each embedding. Call these data a first-level row when \(\eta<b_\tau\) at every embedding.
\end{paramdefinition}

For the two ramified auxiliary primes \(p=q,d\), \eqref{eq:host-auxiliary-completion} gives \(K=K_0(\sqrt{-p})\) with \(p\equiv3\pmod4\). The quadratic Gauss sum gives
\[
 K=K_0(\sqrt{-p})\subseteq K_0(\mu_p),\qquad [K:K_0]=2 .
\]
At an unramified place \(K=K_0\). A finite coefficient extension contains all the indicated embedding images. All actual jumps lie in \([-h,h]\); clipping an empty outer interval or the central interval outside this range therefore preserves support and separation. The normalized endpoints are
\[
 \eta' = \min(\eta,h),\qquad b'_\tau=\min(b_\tau,h+1),
 \qquad \eta'<b'_\tau\le h+1 .
\]
Use the original symbols for these normalized values. Then \(H=h+S+1\) satisfies \(b_\tau+S\le H\). Strict exactness of the crystalline filtration puts the target jumps in the same support.

We describe the integral comparison used here. At full level one
there are \(n\) components indexed cyclically by the embeddings of
\(K_0\). For the tame group \(G\), one has
\[
 |G|=s=(p-1)/e,\qquad O[[\pi]]^G=O[[X]],\qquad
 X\in\pi^s O[[\pi]]^\times,\qquad
 \frac{1}{s}\sum_{g\in G}g\in O[G].
\]
The last element is the integral projection to its invariants.

Let \(E=E_c(X)\)
be the degree-\(e\) polynomial associated to the first cyclotomic
factor on component \(c\), up to a unit. Its roots are the images of
\(X\) at the first cyclotomic points. They belong to the maximal
ideal of \(O\). We use a generator \(\gamma\) of the principal
cyclotomic direction, \(u=\chi_{\rm cyc}(\gamma)\), \(v(u-1)=1\).
On the invariant ring
\[
 \gamma X=u^s X\bmod X^2
\]
and
\(\varphi X\) is a unit times \(X E^s\), using the target component
of Frobenius.

\(N_c(T)\) is the tame-invariant Wach lattice and \(P_c(T)\)
is the invariant Frobenius span from the preceding component. Both
are free over \(A\), with
\[
             E^H P_c\subset N_c\subset E^{-H}P_c .
\]
The action of \(\gamma\) on \(N_c/XN_c\) is trivial. After inverting
\(p\) this quotient gives the corresponding crystalline constants.
Frobenius identifies the constants from \(c-1\) with those used
to evaluate \(P_c\).

\begin{lemma}[Filtration test on bounded numerators]
On each labelled first-level component with logarithmic parameter
\(t\), write \(a\) for the numerator of the element \(E^{-H}a\in E^{-H}P_c\). The test is
\[
 E^{-H}a\in N_c\quad\Longleftrightarrow\quad [t^i]a\in {\rm Fil}^{H-i}_\tau\quad\text{for every first-level root and }i\ge0.
\]
\end{lemma}
\begin{proof}
At a primitive root use a logarithmic parameter
\(t\) on which \(\gamma t=ut\). On bounded-pole Frobenius coordinates,
membership in \(N_c\) is tested by requiring the coefficient of
\(t^i\) to lie in \({\rm Fil}^{-i}_\tau\) at every point of the
component. In numerator coordinates the same condition for a vector \(z=E^{-H}a\) becomes
\[
 [t^i]a\in{\rm Fil}^{H-i}_\tau\qquad(\text{each degree }i).
\]
Indeed \(E\) has a simple zero, and multiplication by a power-series unit preserves the coefficient tests for a decreasing filtration. The evaluation identifies the rational constant
space with the labeled de Rham space using the crystalline comparison;
lifts of the same constant give the same constant value, and the
kernel on the preceding integral lattice is \(XN_{c-1}\).
One uses each embedding over \(c\), including both at ramification
two.

In an adapted de Rham basis of jumps \(a_j\), the completed
filtered lattice is generated by \(t^{-a_j}d_j\), which gives
the coefficient test. Its integral use on already bounded
numerators follows by division by the distinguished factors of
\(E\): if a remainder vanishes at the full set of roots with the
specified multiplicities over the coefficient field, the integral
monic remainder is zero. The height inclusions supply the initial denominator.
\label{cor:row-metric-1}
\end{proof}

We use the integral Wach comparison,
étale exactness, and the finite Herr class criterion for these
lattices. The finite criterion sends a tame-invariant restricted
cocycle pair \((w,z)\), with \(w\in E^{-H}P_c\) and \(z\) sufficiently
deep in \(N\) as below, to a finite class, naturally for \(F\).
The extension comparison sends the connecting pair of a lift of the
unit in a crystalline extension to its integral extension class
up to the fixed unit of the trivial frame.
\label{cor:row-metric-2}

Put
\[
 a=\lceil 2H/s\rceil,\quad M=\lceil4H/s\rceil,\quad
 A_i=i+v_p(i!),\quad \rho=A_{2a-1},\quad \sigma=A_{M-1}.
\]
\begin{proposition}[Integral comparison at the first cyclotomic level]
\label{prop:one-level-integral-row-maps}
For the local field, signed heights and maps in Definition~\ref{hyp:first-level-row},
the map on tame-invariant Wach lattices satisfies
\[
 p^\rho N_c(T_2)\subseteq F(N_c(T_1)).
\]
A first coordinate invariant modulo \(XN_c\) can be multiplied
by \(p^\sigma\) and corrected to become invariant modulo
\(X^MN_c\). It then has a
unique second coordinate in \(X^MN\). The two filtered coordinate
maps \(\alpha,\beta\) have composite \(p^{2\sigma}\) modulo the
specified error modules and preserve the tests modulo \(X^j\)
for \(1\le j\le M\). The same exponents work on the source,
target and unit extension.
\end{proposition}

\begin{proof}
\proofstep{From the tame average to a constant cokernel annihilator.}
First \(p^\rho N_c(T_2)\subset F N_c(T_1)\). For the initial parameter
bound use the ordinary Wach cokernel computation of
Section~\ref{subsec:ramified-lift} on the full cyclotomic field at level
one. Under the étale and Wach identification after restriction to
\(K_0(\mu_p)\), the underlying full-level map is the original map
with the tame action forgotten. The two height bounds of size \(H\)
give \(\pi^{2H}\) in its cokernel annihilator by extracting successive
Frobenius coordinates of a polynomial in \(1+\pi\).

If \(y\) belongs
to the tame-invariant target, \(X^a y\) therefore has a preimage in
the full-level component, since \(sa\ge 2H\). If \(Fz=X^a y\) and \(T_{\rm tame}\) is the tame group, then
\[
 z_0=|T_{\rm tame}|^{-1}\sum_{g\in T_{\rm tame}}gz,\qquad
 Fz_0=|T_{\rm tame}|^{-1}\sum_{g\in T_{\rm tame}}g(X^ay)=X^ay.
\]
The group order is a unit, so this gives an integral invariant preimage. This proves
\(X^a N_c(T_2)\subset F N_c(T_1)\).

On \(X^iN_c/X^{i+1}N_c\) the operator \(\gamma\) is the scalar
\(u^{si}\). This is true on the corresponding filtration of the
image as well, by equivariance. On \(X^ay\) for \(y\in N_c(T_2)\),
apply the two polynomials with roots \(u^{si}\) for \(0\le i<a\)
and \(a\le i<2a\). Both results lie in \(X^aFN_c(T_1)\), by the
image filtration for the first and the target filtration for the
second.

Put \(r_i=u^{si}\) and
\[
 L_i(Z)=\prod_{j\ne i}(Z-r_j)/(r_i-r_j)
\]
on the \(2a\) nodes.
The sum of \(L_i\) for \(i<a\) vanishes on the \(a\) upper
nodes, and that for \(i\ge a\) vanishes on the lower nodes.
Also \(\sum_i L_i=1\). The bound below gives \(p^\rho L_i\in O[Z]\),
so monic division by the products of the opposite roots
gives an identity
\[
         p^\rho=A(Z)\prod_{i<a}(Z-r_i)
                    +B(Z)\prod_{i\ge a}(Z-r_i),
                    \qquad A,B\in O[Z].
\]
Applying it at \(\gamma\) proves the claimed inclusion.

For distinct nodes \(v_p(r_i-r_j)=1+v_p(i-j)\), so the
denominator of the \(i\)-th interpolation polynomial costs
\[
 \sum_{j\ne i}v_p(r_i-r_j)
 =2a-1+v_p\bigl(i!(2a-1-i)!\bigr)
 \le 2a-1+v_p\bigl((2a-1)!\bigr)\le\rho .
\]
The first inequality uses the integrality of the binomial coefficient.
Cancel \(X^a\) to obtain the assertion. The same computation applies
to the surjection from an extension by \(T_2\) to the unit object,
whose extra jump is zero.

\proofstep{Raising the precision and completing the Herr pair.}
A second use of these scalars lifts invariance. If
\(w\in E^{-H}P_c\) satisfies \((\gamma-1)w\in XN_c\), one can
correct \(p^\sigma w\) by an element of \(XN_c\) so that its
difference under \(\gamma\) belongs to \(X^MN_c\). If \(R_i=(\gamma-1)w_i\in X^iN_c\), set
\(w_{i+1}=(u^{si}-1)w_i-R_i\). The scalar action on the \(i\)-th
graded piece gives the next precision:
\[
 (\gamma-1)w_{i+1}=(u^{si}-\gamma)R_i\in X^{i+1}N_c .
\]
The product over \(1\le i<M\) is a unit times \(p^\sigma\).
Also an element of \(XN_c\) invariant to order \(M\) already lies in
\(X^MN_c\), by its first nonzero graded coefficient and
\(u^{si}\ne1\).

On \(X^MN\) the operator \(1-\varphi\) is invertible
by its geometric series. Indeed \(sM-H\ge1\) leaves a factor \(E\)
after applying Frobenius and clearing its height denominator, so
iteration goes to zero in the complete \((p,\pi)\)-adic topology.
Consequently a first Herr coordinate invariant to order \(M\)
has a unique second coordinate in \(X^MN\) solving
\[
 (\varphi-1)z=(\gamma-1)w.
\]
Changing its first coordinate by
\(X^MN\) changes the pair by a coboundary, by the same inverse.

\proofstep{Transporting the horizontal bases by Frobenius.}
For later use we describe the coordinate maps together with their
precision. Let \(L=N_{c-1}/XN_{c-1}\). In an integral basis of
\(N_{c-1}\) lifting a basis of \(L\), construct a matrix \(U(X)\)
with \(U(0)=1\) so the corrected basis is \(\gamma\)-invariant
modulo \(X^M\) over the fraction field.

If the residual first occurs
in degree \(i\), the coefficient equation divides that residual by
\(1-u^{si}\). Thus the degree \(i\) coefficients of \(U\) have
denominator at most \(p^{A_i}\) for \(i<M\). The inverse has the
same bound, since
\[
 \sum A_{i_j}\le A_{\sum i_j}
\]
by factorial
divisibility. Frobenius transport of \(p^\sigma U\) and
\(p^\sigma U^{-1}\) furnishes two integral maps between the model
\(E^{-H} A\otimes_O L\) and \(E^{-H}P_c\), modulo
\(X^M E^{3H}\) times the respective undivided modules. Their composites are \(p^{2\sigma}\) in these quotients because
\[
 \frac{\varphi(X)^M}{E^H}\in X^M E^{sM-H}A
                         \subseteq X^M E^{3H}A
 \qquad(sM\ge4H).
\]
They commute in the quotients with \(\gamma\) and lattice maps.
For the latter property the map in horizontal bases agrees through
degree \(M-1\) with its constant matrix, by equivariance and
\(u^{si}\ne1\). Denote the forward and backward maps by \(\alpha,\beta\).

\proofstep{The filtration tests at all first-level roots.}
In the first-level expansion the corrected basis specializes to its
constant through orders below \(sM\) before division by \(E^H\).
Indeed the pulled-back invariant equation forces all intermediate
positive coefficients to vanish, since the degree \(i\) cyclotomic
eigenvalue is \(u^i\ne1\). The precision and jump-height bounds give
\[
 i\ge sM\ \Longrightarrow\ i\ge4H\ \Longrightarrow\ i\ge2H.
\]
All first-level filtration tests for numerator coefficients in those degrees are automatic. Hence both maps carry the filtered model lattice,
denoted \(N^{\rm mod}\), into \(N_c\) and back, respectively, modulo
the indicated error modules. They preserve the tests for membership
in \(X^jN\) for \(1\le j\le M\) in the forward or backward direction
from a member, by linearity in the invariant coefficient ring and
the inclusion of the error in \(X^MN\). In particular an invariant
modulo \(XN\) has an invariant partner. These constructions apply
with the same \(H,M,\rho,\sigma\) to the source, target and unit
extension.
\label{cor:row-metric-3}
\end{proof}
